% Minimal retained preamble for lane-p2 cards taken from
% sources/logic-fl02-conservation-of-ramsey-s-theorem-f/pinned-original-source.tex
% (Conservation of Ramsey's theorem for pairs, arXiv:2402.11616v3).
%
% Copied verbatim from the author's own file:
%   line 2      : \documentclass{article}
%   lines 16-25 : the author's theorem declarations
%   lines 46-88 : the author's macro definitions for the reverse-mathematics principles
%   line 91     : \card
%   lines 104-107: the natural sums and products, which rely on mathabx (source line 103)
%   line 116    : \qt
%
% Omitted author preamble lines and why (nothing was renamed or redefined):
%   line 3 (graphicx), 7 (url), 9 (cancel): no figure, url or cancel use in the retained spans;
%   line 6 (cleveref): the assembler adds its own cleveref machinery, a second load clashes;
%   line 8 (\usepackage[utf8x]{inputenc}): inputenc with the utf8x option fights the forced
%       xelatex engine;
%   lines 28-37 (the rep@theorem machinery and \newreptheorem{maintheorem}): the retained spans
%       use no rep-theorem environment;
%   lines 40-43 (xcolor and the three author comment macros): not used by the retained spans;
%   lines 93-99 (the LS2 font encoding block for \smashtimes based on stix): the host TeX Live has
%       no stix directory and the author himself never uses the result -- line 105 maps \ntimes to
%       \dottimes, and line 110 shows \smashtimes already commented out in the source.
% mathabx is retained because \dotplus, \dottimes and \dotdiv come from it.
\documentclass{article}
\usepackage{amsmath,amssymb,amsthm}
\usepackage{mathabx}
\usepackage{cancel}
\usepackage{cleveref}

\newtheorem{theorem}{Theorem}
\numberwithin{theorem}{section}
\newtheorem{maintheorem}[theorem]{Main Theorem}
\newtheorem{lemma}[theorem]{Lemma}
\newtheorem{question}[theorem]{Question}
\newtheorem{proposition}[theorem]{Proposition}
\newtheorem{statement}[theorem]{Statement}
\newtheorem{remark}[theorem]{Remark}
\newtheorem{definition}[theorem]{Definition}
\newtheorem{corollary}[theorem]{Corollary}

\newcommand{\RCA}{\mathsf{RCA}}
\newcommand{\WKL}{\mathsf{WKL}}
\newcommand{\SRT}{\mathsf{SRT}}
\newcommand{\EM}{\mathsf{EM}}
\newcommand{\SEM}{\mathsf{SEM}}
\newcommand{\ADS}{\mathsf{ADS}}
\newcommand{\SCAC}{\mathsf{SCAC}}
\newcommand{\SADS}{\mathsf{SADS}}
\newcommand{\COH}{\mathsf{COH}}
\newcommand{\ACA}{\mathsf{ACA}}
\newcommand{\qvdash}{\operatorname{{?}{\vdash}}}
\newcommand{\nqvdash}{\operatorname{{?}{\nvdash}}}
\newcommand{\BSig}{\mathsf{B}\Sigma^0}
\newcommand{\ISig}{\mathsf{I}\Sigma^0}
\newcommand{\LSig}{\mathsf{L}\Sigma^0}
\newcommand{\BPi}{\mathsf{B}\Pi^0}
\newcommand{\IPi}{\mathsf{I}\Pi^0}
\newcommand{\LPi}{\mathsf{L}\Pi^0}
\newcommand{\IDel}{\mathsf{I}\Delta^0}
\newcommand{\BCDe}{\mathsf{BC}\Delta^0}
\newcommand{\WF}{\mathsf{WF}}
\newcommand{\RT}{\mathsf{RT}}
\newcommand{\BME}{\mathsf{BME}}
\newcommand{\TT}{\mathsf{TT}}
\renewcommand{\O}{\mathcal{O}}
\renewcommand{\L}{\mathcal{L}}
\newcommand{\C}{\mathcal{C}}
\newcommand{\R}{\mathcal{R}}
\renewcommand{\S}{\mathcal{S}}
\newcommand{\T}{\mathcal{T}}
\renewcommand{\P}{\mathcal{P}}
\newcommand{\M}{\mathcal{M}}
\newcommand{\Nc}{\mathcal{N}}
\renewcommand{\P}{\mathcal{P}}
\newcommand{\Q}{\mathcal{Q}}
\newcommand{\NN}{\mathbb{N}}
\newcommand{\PP}{\mathbb{P}}
\newcommand{\uh}[0]{{\upharpoonright}}
\newcommand{\card}{\operatorname{card}}
\newcommand{\nplus}{\dotplus}
\newcommand{\ntimes}{\dottimes}
\newcommand{\nminus}{\dotdiv}
\newcommand{\nsum}{\dot{\sum}}
\def\qt#1{``#1''}
\title{An infinite maximal tree would give an infinite descending sequence of ordinals}
\author{Quentin Le Houérou, Ludovic Levy Patey, Keita Yokoyama}
\date{arxiv:2402.11616v3; CC BY 4.0}
\begin{document}
\maketitle
\noindent\textit{Editorial note.} The counted unit is the article's own Lemma lem:sads-wf-forcing-wf-jump at source lines 1250-1259 together with the proof environment that begins two lines after the statement ends (1261-1292); the pairing is the ordinary adjacent one. This label uses the author's plain form rather than a typed one, and the gate recognises both.\par\medskip
\noindent\textit{Editorial note.} Eight context spans are retained: the Shore-blocking lemma at 426-431, the definitions of a condition, of a split pair and of the two forcing relations at 1114, 1117-1119 and 1121-1127, the definitions of trees and satisfaction at 1169-1183 and 1186-1192, and the statements of the two SADS tree lemmas at 1194-1196 and 1206-1215. All are prerequisites that the counted statement and proof use; their own proofs are not retained and they are counted by sibling cards of this lane.\par\medskip
\noindent\textit{Editorial note.} The retained preamble reproduces verbatim the author's document class at source line 2, his theorem declarations at source lines 16-25, his macro definitions at source lines 46-88 and 91, and his natural-operation macros at source lines 104-107, and it retains mathabx and cleveref; the stix-dependent LS2 block and the inputenc load with the utf8x option are not retained. mathabx supplies the natural sum and product symbols used in the rank definition.\par\medskip
\noindent\textit{Editorial note.} No retained line contains a citation, so the card carries no bibliography.\par\medskip
\noindent\textit{Editorial note.} The two remarks that follow the proof at source lines 1294-1302 are not counted here and are not used as premises.\par\medskip
\noindent\textit{Allowed standard prerequisites (editorial).}\begin{itemize}
\item The two formulas Phi\_k and Psi\_k and the Shore-blocking functional. Phi\_k(G,x) states that there are indices below k, increasing in number x, at which the functional converges on the generic; Psi\_k(G,x) states that the Shore-blocking functional Gamma\_k converges at x, where Gamma\_k is the functional of the Shore-blocking lemma for the ordinal alpha; the author's own definitions are retained at source lines 1250-1251 and the underlying lemma at 426-431.. They are the two formulas out of which the counted lemma lets the pair be chosen, and the quantities f\_0 and f\_1 in the proof are built from their values.
\item Maximal trees of split pairs and their finite-tree forcing consequence. For every pair of Sigma\textasciicircum{}0\_1 formulas and every condition there is a maximal tree of split pairs satisfying the pair, and if that tree is M-finite then some extension forces the two formulas at two indices and forces the negation of one of the next instances; retained with this card at source lines 1194-1196 and 1206-1215.. The proof obtains the maximal tree from the first statement and concludes with the second once it proves that the tree is M-finite, which is the only remaining step.
\item Ordinal-valued rank of a finite tree and well-foundedness of the supremum. To each finite tree one assigns the ordinal obtained by summing, over the nodes, the ordinal omega raised to the sum of the two values f\_0 and f\_1 at the node, multiplied by the number of missing immediate successors; the hypothesis WF of the supremum of omega to the natural product with omega below alpha forbids an M-infinite decreasing sequence of such ordinals; the article uses these ordinals in the proof.. The proof defines exactly such a rank zeta\_s for each approximating tree and shows that the ranks strictly decrease, which contradicts the well-foundedness hypothesis.
\end{itemize}\medskip
\begin{lemma}\label[lemma]{lem:shore-blocking-wf}
Consider a model $\M = (M,S)$ and $\alpha \in M$ an ordinal less than $\epsilon_0$.
For every~$k \in M$, there is a Turing functional $\Gamma_k$ such that, letting~$\beta \in \alpha$ be the largest ordinal less than $\alpha$ of index at most~$k$, for every~$X \in 2^M$ such that $\M[X] \models \RCA_0$, if there is some~$e < k$ such that $\Phi^X_e$ is an $M$-infinite decreasing sequence of elements smaller than~$\beta$, then $\Gamma^X_k$ is an $M$-infinite decreasing sequence of elements smaller than~$\beta \ntimes k$.

Moreover, an index of~$\Gamma_k$ can be found computably in~$k$.
\end{lemma}

A condition is a pair $(\sigma, \tau)$ with $\sigma$ an $M$-finite $L$-ascending sequence included in $U$ and $\tau$ an $M$-finite $L$-descending sequence included in $V$. We have $(\sigma_2, \tau_2) \leq (\sigma_1, \tau_1)$ if $\sigma_2 \succeq \sigma_1$ and $\tau_2 \succeq \tau_1$.

\begin{definition}
A \emph{split pair} is a pair $(\sigma, \tau)$ of strings such that $\sigma$ is an $M$-finite $L$-ascending sequence, $\tau$ is an $M$-finite $L$-descending sequence, and $\max_L \sigma <_L \min_L \tau$.
\end{definition}

\begin{definition}
Let $(\sigma_0, \sigma_1)$ be a condition, $e, x \in M$ and $i < 2$, we write:
\begin{itemize}
    \item $(\sigma_0, \sigma_1) \Vdash_i \Phi_e^{G \oplus Y}(x) \downarrow$ if $\Phi_e^{\sigma_i \oplus Y}(x) \downarrow$. 
    \item $(\sigma_0, \sigma_1) \Vdash_i \Phi_e^{G \oplus Y}(x) \uparrow$ if $\forall \sigma'_i \succeq \sigma_i$ such that $(\sigma'_i, \sigma_{1-i})$ is a split pair, $\Phi_e^{\sigma'_i \oplus Y}(x)\uparrow$.
\end{itemize}   
\end{definition}

\begin{definition}
Fix a condition $(\sigma, \tau)$.

A \emph{$(\sigma, \tau)$-tree} is a binary tree $T$ labelled by a family of split pairs $(\sigma_{\rho}, \tau_{\rho})_{\rho \in T}$ such that for all $\rho \in T$:
\begin{itemize}
    \item $\sigma \preceq \sigma_{\rho}$ and $\tau \preceq \tau_{\rho}$ ;
    \item For every $\mu \in T$, if $\rho \cdot 0 \preceq \mu$ then $\sigma_{\rho} \preceq \sigma_{\mu}$ and if $\rho \cdot 1 \preceq \mu$ then $\tau_{\rho} \preceq \tau_{\mu}$ ;
\end{itemize}
%    For a node $\rho \in T$, we define $\pi(\rho)$ to be equal to $(\sigma_{\rho \uh i_0}, \tau_{\rho \uh j_0})$ where $i_0$ is the biggest index $i < |\rho|$ such that $\rho(i) = 0$ (if no such index exists then we take $\sigma$ instead of $\sigma_{\rho \uh i_0}$ in the definition of $\pi(\rho)$) and $j_0$ is the biggest index $j < |\rho|$ such that $\rho(j) = 1$ (simillarly, if no such index exists, take $\tau$ instead of $\tau_{\rho \uh j_0}$). 

%    A node is said to be \emph{valid} if $\pi(\rho)$ is a condition.
     %A node $\rho$ is said to be \emph{valid} if the biggest index $i$ such that $\rho(i) = 0$ satisfies $\sigma_{\rho \uh i} \subseteq U$ (if no such index exists this condition is verified) and if the biggest index $j$ such that $\rho(i) = 1$ satisfies $\tau_{\rho \uh i} \subseteq V$ (if no such index exists this condition is verified). Being valid is a $\Delta_2^0(Y)$ condition.

     A node $\rho$ is said to be \emph{valid} if for every $\mu \in T$ such that $\mu \cdot 0 \preceq \rho$, $\sigma_{\mu} \subseteq U$ and for every $\mu \in T$ such that $\mu \cdot 1 \preceq \rho$, $\tau_{\mu} \subseteq V$. Being valid is a $\Delta^0_2(Y)$ condition.
\end{definition}

\begin{definition}
    Let $T$ a $(\sigma,\tau)$-tree 
    \begin{itemize}
        \item A node $(\sigma_{\rho}, \tau_{\rho})$ in $T$ \emph{satisfies} a pair of $\Sigma^0_1(Y)$-formulas $(\varphi_0(G_0, x), \varphi_1(G_1,x))$ if $\varphi_0(\sigma, |\rho| - \card \rho)$ and $\varphi_1(\tau, \card \rho)$ holds (where $\card \rho$ is the cardinal of $\rho$ seen as a set).
        \item $T$ \emph{satisfies} $(\varphi_0(G_0, x), \varphi_1(G_1,x))$ if every node of $T$ satisfies it.
    \end{itemize}
\end{definition}

\begin{lemma}\label[lemma]{lem:sads-wf-exists-maximal-tree}
For every pair $(\varphi_0(G_0, x), \varphi_1(G_1,x))$ of $\Sigma^0_1(Y)$ formulas and every condition $(\sigma, \tau)$, there is a maximal $\Sigma^0_1$ $(\sigma, \tau)$-tree~$T$ satisfying it.
\end{lemma}

\begin{lemma}\label[lemma]{lem:sads-wf-finite-tree-forcing}
Let $(\sigma,\tau)$ be a condition, $(\varphi_0(G_0, x), \varphi_1(G_1,x))$ be a pair of $\Sigma^0_1(Y)$ formulas, and $T$ be an $M$-finite maximal $(\sigma, \tau)$-tree satisfying $(\varphi_0(G_0, x), \varphi_1(G_1,x))$.
Then there is an extension $(\hat \sigma, \hat \tau) \leq (\sigma, \tau)$ and some~$p,q \in M \cup \{-1\}$ such that the following holds:
\begin{itemize}
    \item $(\hat \sigma, \hat \tau) \Vdash_0 \varphi_0(G_0, p)$ and $(\hat \sigma, \hat \tau) \Vdash_1 \varphi_1(G_1,q)$
    \item $(\hat \sigma, \hat \tau) \Vdash_0 \neg \varphi_0(G_1,p+1)$ or $(\hat \sigma, \hat \tau) \Vdash_1 \neg \varphi_1(G_1,q+1)$.
\end{itemize}
With the convention that  $(\hat \sigma, \hat \tau) \Vdash_i \varphi_i(G^i, -1)$ always holds.

\end{lemma}

\begin{lemma}\label{lem:sads-wf-forcing-wf-jump}
Let $k \in M$, let $\Phi_k(G,x)$ be the formula: $\exists e_0 < e_1 < \dots < e_{x} < k$ such that $\Phi_{e_i}^{G \oplus Y}(e_i) \downarrow$ for all $i < x$, let $\Psi_k(G,x)$ be the formula: $\Gamma^{G \oplus Y}_k(x) \downarrow$ (with $\Gamma_k$ the functional of \Cref{lem:shore-blocking-wf} for the order $\alpha$).

For any combination $\phi_0(G,x), \phi_1(G,x) \in \{\Phi_k(G,x), \Psi_k(G,x)\}$ and every condition $(\sigma, \tau)$, there exists an extension  $(\hat \sigma, \hat \tau) \leq (\sigma, \tau)$ and some~$p,q \in M \cup \{-1\}$ such that the following holds:
\begin{itemize}
    \item $(\hat \sigma, \hat \tau) \Vdash_0 \varphi_0(G_0, p)$ and $(\hat \sigma, \hat \tau) \Vdash_1 \varphi_1(G_1,q)$
    \item $(\hat \sigma, \hat \tau) \Vdash_0 \neg \varphi_0(G_1,p+1)$ or $(\hat \sigma, \hat \tau) \Vdash_1 \neg \varphi_1(G_1,q+1)$.
\end{itemize}
With the convention that $(\hat \sigma, \hat \tau) \Vdash_i \varphi_i(G_i, -1)$ always holds.
\end{lemma}

\begin{proof}
By \Cref{lem:sads-wf-exists-maximal-tree} there exists a maximal $\Sigma^0_1$ $(\sigma, \tau)$-tree~$T$ satisfying $$(\varphi_0(G, x), \varphi_1(G,x)).$$ To apply \Cref{lem:sads-wf-finite-tree-forcing} and conclude, we just need to prove the $M$-finiteness of~$T$. 

By twisting $\Gamma_k$, we can assume that for every~$e, n \in M$, if $\Gamma^{\sigma \oplus Y}_k(n)\downarrow$, then
\begin{itemize}
    \item $\Gamma^{\sigma \oplus Y}_k(m)\downarrow$ for every~$m < n$ ;
    \item $\Gamma^{\sigma \oplus Y}_k(0), \Gamma^{\sigma \oplus Y}_k(1), \dots, \Gamma^{\sigma \oplus Y}_k(m)$ is a strictly $<_{\epsilon_0}$-decreasing sequence of elements smaller than~$\beta \ntimes k$.
\end{itemize}

Assume by contradiction that $T$ is $M$-infinite and let $(T_s)_{s \in M}$ be an increasing sequence of computable tree approaching $T$. Without loss of generality we can assume that $T_s \neq T_{s+1}$ for all $s \in M$.

For every $s \in M$ and $\rho$ node of $T_s$ indexed by $(\sigma_{\rho}, \tau_{\rho})$ with set of immediate successors $S_\rho$ ($\card S_\rho \leq 2$ since $T_s$ is a binary tree), define inductively :
$$\zeta_s(\rho) = \omega^{f_0(\rho) \nplus f_1(\rho)}(2 - \card S_\rho)\ \nplus \nsum_{\mu \in S_\rho} \zeta_s(\mu)
$$
Where $f_0(\rho) = \Gamma_k^{\sigma_{\rho} \oplus Y}(|\rho| - \card \rho)$ if $\varphi_0(G,x) = \Psi_k(G,x)$ and $f_0(\rho) = k - (|\rho| - \card \rho)$ if $\varphi_0(G,x) = \Phi_k(G,x)$. And $f_1(\rho) = \Gamma_k^{\tau_{\rho} \oplus Y}(\card \rho)$ if $\varphi_1(G,x) = \Psi_k(G,x)$ and $f_1(\rho) = k - \card \rho$ if $\varphi_1(G,x) = \Phi_k(G,x)$.

And let $\zeta_s = \zeta_s(\epsilon)$. \\

\textbf{Claim}. \emph{The sequence $\zeta_s$ is an $M$-infinitely $<_{\epsilon_0}$-decreasing sequence of ordinals less than $\omega^{\beta \times \omega}$}

Let $s \in M$, since $T_{s+1} \neq T_s$, there is a $\rho \in T_s$ an $i < 2$ such that $T_{s+1} = T_s \cup \rho \cdot i$. Therefore:

$$\zeta_{s}(\rho) \nplus (\omega^{f_0(\rho \cdot i) + f_1(\rho \cdot i)} \ntimes 2) = \zeta_{s+1}(\rho) \nplus \omega^{f_0(\rho) \nplus f_1(\rho)}$$

But for every possible values of $\varphi_0$ and $\varphi_1$, we have $f_0(\rho) > f_0(\rho \cdot 0)$, $f_0(\rho) = f_0(\rho \cdot 1)$, $f_1(\rho) = f_1(\rho \cdot 0)$ and $f_1(\rho) > f_1(\rho \cdot 1)$. So for every possible value of $i$, we have 
$$\omega^{f_0(\rho) + f_1(\rho)} > \omega^{f_0(\rho \cdot i) + f_1(\rho \cdot i)} \times 2$$
and therefore $\zeta_{s+1}(\rho) < \zeta_s(\rho)$ which gives us $\zeta_{s+1} < \zeta_s$ (as $\rho \cdot i$ is the only difference between $T_s$ and $T_{s+1}$). \\

Therefore, for all $s \in M$, $\zeta_s \leq_{\epsilon_0} \zeta_0 \leq_{\epsilon_0} \omega^{\beta \times \omega}$. \\

This contradicts $\WF(\sup_{\beta < \alpha} \omega^{\beta \times \omega})$, so $T$ is $M$-finite and the result follows.
\end{proof}

\end{document}
